\section{Critiques of simple production functions \label{prodfuncts}}

Jane Jacobs' arguments in particular present problems for the straightforward  `black box' model. Where the point of those models is that an in-out relationship can be defined - and only altered through some technological change - Jacobs' description of production paints a picture of economic actors stumbling upon new combinations of goods in a way very reminiscent of evolution. 

Perhaps the most challenging productive factor to theorise is knowledge. An example from Krugman illustrates the sort of interactions that can take place. He argues that as energy became cheaper and more accessible, it could substitute for skilled labour. The result was a replacement of that labour by cheap energy sources - and an increase in the footloose mobility of workers, because pools of local skilled labour were now less necessary. 

\subsection{Criticism}

The kind of production and utility functions described get plenty  of criticism, aimed in one way or another at their intrinsic lack of realism. Ecological economics, for example, sees in the Cobb Douglas function a 'Garden of Eden' world (GR quoted in Daly) where natural resources are optional. Daly sums this up nicely, if somewhat loquaciously. He responds to Solow's line that ``If it is very easy to substitute other factors for natural resources, then there is in principle no `problem'. The world can, in effect, get along without natural resources'' (Solow quoted in \citealt{daly_georgescu-roegen_1997} p.261)

Daly responds that economists seem to believe in ``making a cake with only the cook and his kitchen. We do not need flour, eggs, sugar, etc., nor electricity or natural gas, nor even firewood. If we want a bigger cake, the cook simply stirs faster in a bigger bowl and cooks the empty bowl in a bigger oven that somehow heats itself.'' (ibid.) 

More formally, he quotes Georgescu-Rosen's critique of the Solow-Stiglitz production function - a Cobb Douglas with the addition of natural resources: 

\begin{equation}
Q=K^{a1}R^{a2}L^{a\text{3}}\label{eq:}
\end{equation}

He points out this means that $R^{a2}=Q/K^{a1}L^{a3}$. So, holding labour constant, one could get any quantity of output with as small an amount of $R$, as long as $R>0$ and $K$ is large enough to compensate for that tiny amount.  Solow's reply to this highlights a recurring argument used to defend the abstract nature of many economic models: critics are taking them too literally, and not considering how the models are used. Solow counters with:

\begin{quote}
``We were trying to think about an interesting and important question: how much of a drag on future growth, or even on the sustainability of current production, might be exercised by the limited availability of natural resources and the inputs they provide? ... The role of theory is to explore what logic and simple assumptions can tell us about what data to look for and how to interpret them in connection with the question asked.'' \fixref{(solow p.267/8)}
\end{quote}

Solow goes on to point out that the argument should be about how substitutable renewable and non-renewable resources are, given that the former are likely to be highly capital-intensive. (Ibid.) He clearly does understand the issues; what the critics seem to have done is mistake the economists' models for their actual understanding of the world, rather than tools that aid that understanding. Krugman makes much the same kind of plea. {[}insert details here!{]} A clear example of this also comes from Krugman's work on international trade, discussed above, and cannot be put better than Krugman himself does in the abstract to his original 1979 paper: 

\begin{quotation}
``... trade, and gains from trade, will occur, even between countries with identical tastes, technology, and factor endowments.'' (Krugman 1979 p.469) 
\end{quotation}

Or rather, they will occur in the model that Krugman presents. The model itself cannot guarantee the real world contains the same dynamic; the point is to show that such a dynamic is \emph{possible}, given a particular set of assumptions. 

x

While Solow and Krugman are in large part correct, there are problems  with their case. Most fundamentally, there is a large gap between, say, Krugman's stated awareness of a model's ``silly assumptions'' and the uses to which those models are now being put in policy circles. 

Pursuit of that argument, however, is beyond the scope of the thesis. The more important issue at hand is this: there are implications for agent modelling. How so? It comes back to Daly's claim that standard economic functions offer only 	 (p.262) Consider an agent model of production that has aspirations to examine the complex morphology of linkages seen in reality - Jane Jacobs' `tangled bank'. Agent models are almost always explicitly time-driven; computer cycles are synonymous with timesteps and iteration is central. In such a world, what moves between economic linkages each timestep, if not actual resources? Certainly, many exchanges can be service-related, but in a time-driven model it becomes much harder to conceive of resources in an abstract away. Instead, a more Weberian image presents itself as the most likely:
resources with weights, combined to produce other resources with weights.